\documentclass{article}
\usepackage{iclr2026_conference,times}

% Optional math commands from https://github.com/goodfeli/dlbook_notation.
\input{math_commands.tex}

\usepackage{hyperref}
\usepackage{url}
\usepackage{graphicx}
\usepackage{booktabs}
\usepackage{amsmath}
\usepackage{amssymb}

% ============================================================================
% Title
% ============================================================================
\title{Your Paper Title Here}

% ============================================================================
% Authors
%
% Authors must NOT appear in the submitted version. They should stay hidden
% as long as \iclrfinalcopy below remains commented out. Non-anonymous
% submissions will be desk-rejected.
% ============================================================================
\author{Anonymous Authors \\
Anonymous Affiliation \\
\texttt{anonymous@example.com}
% Replace once accepted by something like:
% Alice Author \\
% Department of Computer Science \\
% Your University \\
% \texttt{alice@univ.edu}
% \And
% Bob Coauthor \\
% \texttt{bob@univ.edu}
}

% Uncomment ONLY for the camera-ready version (after acceptance) — never for
% the submission itself.
% \iclrfinalcopy

\begin{document}
\maketitle

% ============================================================================
\begin{abstract}
% 1 paragraph, ~200 words.
% State the problem, the gap, your contribution, and the headline result.
% Avoid "we propose a novel ..." — name what's specifically new.
Replace this paragraph with your abstract.
\end{abstract}

% ============================================================================
\section{Introduction}
\label{sec:intro}

Motivate the problem in one paragraph. State the gap in prior work in the
next. Then list contributions explicitly:

\begin{itemize}
    \item Contribution 1 — be specific.
    \item Contribution 2.
    \item Contribution 3.
\end{itemize}

Close with a paragraph on paper organization (Section~\ref{sec:related}
reviews \dots; Section~\ref{sec:method} introduces \dots).

% ============================================================================
\section{Related Work}
\label{sec:related}

Group related work by theme, not by venue. For each closest paper, state
explicitly how this work differs. Avoid the "Smith et al.\ also studied X,
but \dots" laundry list — group, compare, contrast.

% ============================================================================
\section{Method}
\label{sec:method}

State assumptions upfront. Define notation on first use:
let $\mX \in \mathbb{R}^{n \times d}$ denote \dots

Use subsections for sub-components. The reader should be able to reproduce
your method from this section alone (modulo hyperparameters, which go in
the experiments section or appendix).

% ============================================================================
\section{Experiments}
\label{sec:exp}

\subsection{Setup}
Datasets, baselines, metrics, hardware, hyperparameters, random seeds.

\subsection{Main Results}
% Example table — replace with yours.
\begin{table}[t]
\centering
\caption{Main results. Best in \textbf{bold}.}
\label{tab:main}
\begin{tabular}{lcc}
\toprule
Method & Metric A $\uparrow$ & Metric B $\downarrow$ \\
\midrule
Baseline 1 & 65.4 & 0.42 \\
Baseline 2 & 67.1 & 0.39 \\
\textbf{Ours} & \textbf{72.3} & \textbf{0.31} \\
\bottomrule
\end{tabular}
\end{table}

\subsection{Ablations}
One ablation per contribution claim in the introduction.

% Example figure stub — uncomment + drop image into figures/
% \begin{figure}[t]
% \centering
% \includegraphics[width=0.7\linewidth]{figures/teaser.pdf}
% \caption{One-line caption that stands alone.}
% \label{fig:teaser}
% \end{figure}

% ============================================================================
\section{Discussion}
\label{sec:discussion}

What did we learn? Where does the method break? What is the right way to
interpret the gains? Honest limitations belong here — reviewers reward
candor.

% ============================================================================
\section{Conclusion}
\label{sec:conclusion}

Restate the contribution (no inflation). Future work should be specific,
not "we plan to extend this to more domains."

% ============================================================================
% Bibliography
% ============================================================================
\bibliographystyle{iclr2026_conference}
\bibliography{references}

% ============================================================================
% Appendix (uncomment if needed)
% ============================================================================
% \appendix
% \section{Additional Experiments}
% \label{app:additional}

\end{document}
